\section{The healthy-field comparison at the middle scale}
\label{sec:marked-healthy-cumulant-launch}

The Gaussian comparison must approximate the actual local coefficient
on each healthy configuration. We prove this by constructing one closed
real curvature cochain from that configuration, and by comparing the
actual polynomial transport with the transport of its covariance atlas.
Both comparisons are simultaneous for every subset of the insertion word.

\subsection{A closed cochain from a healthy gauge field}

Work on a lattice of integer side $N_{\rm per}>0$, with a patch
of integer radius $B\ge0$ about $o$, and assume $2(B+1)<N_{\rm per}$.
Let the field be healthy at threshold
$\theta\ge0$ on that patch. Put it in the axial gauge based at $o$.
The path construction expresses every link in the room as a product
of at most $4B$ conjugated plaquette holonomies. The chord of a
product is at most the sum of the chords, so each such link has chord
at most $4B\theta$. Let $c>0$ be the actual scale normalization.
Define a Lie-algebra-valued one-cochain $A$ by taking $c^{-1}$
times the curvature coordinate of this gauge-fixed link inside the
room, and zero outside it. Define the two-cochain $X=dA$ on the
whole torus, and take its real colour coordinates. Then
\begin{equation}\label{eq:healthy-closed-cochain}
           dX=0,\qquad \|X\|_\infty\le16B\theta/c.
\end{equation}
The first assertion is the exact discrete identity $d^2=0$, also
across the artificial boundary of the room. The second follows by
summing the four link bounds in each plaquette curl.

Let $Y$ be $c^{-1}$ times the actual curvature of the gauge-fixed
plaquette. For a plaquette whose base is at distance at most $B-2$
from $o$, all four edges lie inside the room. The curvature map
satisfies the product estimate
\[
 \|\mathcal C(ab)-\mathcal C(a)-\mathcal C(b)\|
       \le c_G\operatorname{chord}(a)\operatorname{chord}(b),
 \qquad \mathcal C(a^{-1})=-\mathcal C(a).
\]
Applying it successively to the four factors of the plaquette,
with their inverse orientations, costs $1+2+3=6$ products of
link bounds. Hence
\begin{equation}\label{eq:healthy-curvature-linear-error}
     |Y_\nu-X_\nu|\le6c_G(4B\theta)^2/c
\end{equation}
for every colour coordinate in the interior stencil. This is an
estimate for the same gauged field, without sampling a new field.
Translations and signed coordinate changes transport the construction
and preserve these absolute bounds.

At the selected middle scale the normalization satisfies
$\theta/c\le\tau(t)$, where $t$ is its coupling and
$\tau(t)=3\,\mathrm{hot}(t)$. The upper comparison coupling is
$s$, with $s/2\le t\le s$. Write $H=\exp((\log\log s)^2)$,
taking $s$ large enough that $H\ge1$, and fix $0<\lambda\le1/4$.
The purchased numerical bounds are
\begin{align}\label{eq:healthy-middle-purchases}
 16B\tau(t)&\le H^{\lambda/4},\\
 6c_G(4B)^2\theta\tau(t)&\le\varepsilon,
                  \qquad \varepsilon=H^{-v},\notag
\end{align}
where $v$ is selected before the final coupling threshold.
They are feasible uniformly in this interval: the patch radius and
$\tau$ grow slower than every power of $H$, whereas
$\theta=t^{-1/2}\mathrm{hot}(t)$ decays faster than every negative
power of $H$. For example, taking logarithms compares
$-\tfrac12\log t+o(\log t)$ with any fixed multiple of
$(\log\log s)^2$; $t\ge s/2$ makes this uniform.
Thus \eqref{eq:healthy-closed-cochain} supplies exactly a closed real
background admitted by the covariance comparison, while
\eqref{eq:healthy-curvature-linear-error} is at most $\varepsilon$.

\subsection{Comparing two polynomial transports}

A polynomial step consists of a finite linear substitution followed
by Wick contraction. Suppose the coefficient sum of every substituted
coordinate and every covariance entry is at most $C\ge1$.
For degree at most $D$, put $A_D=[2(D+2)C]^D$.
The Wick recurrence used in
Section~\ref{sec:allarity-cumulant-comparison} applies equally when
its mean is a linear polynomial. Submultiplicativity of coefficient
sums then proves that the step preserves degree and has norm at most
$A_D$ on this finite-degree space. If a second step has the same
bounds and each linear row and covariance entry differs by at most
$\varepsilon\ge0$, the same induction, subtracting products one factor
at a time, bounds the difference operator by $\varepsilon A_D$.
The row error here is a coefficient sum, not an entrywise error
followed by an unbounded sum over sites.

For $m$ successive steps $T_i,T'_i$, expand
\[
 T_m\cdots T_1-T'_m\cdots T'_1
 =\sum_{j=1}^m T'_m\cdots T'_{j+1}
                   (T_j-T'_j)T_{j-1}\cdots T_1.
\]
Thus, for $\deg p\le D$,
\begin{equation}\label{eq:healthy-transport-telescope}
 |T_m\cdots T_1p-T'_m\cdots T'_1p|_1
                         \le |p|_1 m\varepsilon A_D^m.
\end{equation}
We apply this to the actual polynomial history and the history read
from its covariance atlas. Their per-step common bounds and errors
are the same ones used in the marked Gaussian reader estimates.

All degree bounds $D,K,D_*$ and the insertion count $n$ below are
nonnegative integers. Let $p_i$ be the calibrated incoming polynomials,
$i=1,\ldots,n$, with a fixed exponent $z\ge0$,
with $|p_i|_1\le H^z$ and degree at most $K$. Every subset product
has degree at most $nK$ and norm at most $H^{nz}$. Choose the
block length after $A_{nK}$ so that
\[
              2(b+2)\log A_{nK}\le\log L.
\]
For the prescribed number
$m=\lceil(b+2)\log H/\log L\rceil$ of steps, this gives
$A_{nK}^m\le H$ eventually; also $m\le H$ eventually.
Therefore all transported subset moments have norm at most
$R=H^{nz+1}$, and their differences have norm at most
$RH\varepsilon$.

For a nonempty subset $I$, its connected polynomial is
\[
 c_I=\sum_{\pi\in\mathcal P(I)}
       (-1)^{|\pi|-1}(|\pi|-1)!\prod_{B\in\pi}q_B,
\]
where $q_B$ is the transported moment polynomial. There are finitely
many set partitions. Telescoping the product for each partition
bounds the difference of the two connected polynomials by
\begin{equation}\label{eq:healthy-connected-polynomials}
 |c_I-c'_I|_1\le D_n(R)RH\varepsilon,
 \quad D_n(R)=\sum_{\varnothing\ne I\subseteq\{1,\ldots,n\}}
       \sum_{\pi\in\mathcal P(I)}|\pi|!R^{|\pi|-1}.
\end{equation}
Both have degree at most a fixed number $D_*$, and coefficient norm
at most a fixed polynomial $Q_n(R)$ in $R$, obtained by the same
partition expansion. These constants depend on the finite word,
not on the lattice period.

\subsection{Evaluation and the initial connected error}

For a polynomial $p$ of degree at most $D_*$, two real coordinate
vectors bounded by $M\ge1$ on its variables, and coordinate error
at most $\varepsilon$, monomial telescoping gives
\[
 |p(Y)-p(X)|\le |p|_1D_*M^{D_*}\varepsilon.
\]
The sharper power $M^{D_*-1}$ is available for positive degree;
the displayed bound also covers degree zero. Moreover
$|p(X)|\le |p|_1M^{D_*}$.
The healthy-field construction and its numerical purchases bound
both vectors by $M=H$ on the required interior stencil.
Using \eqref{eq:healthy-connected-polynomials} therefore gives
\begin{align*}
 |c_I(Y)-c'_I(X)|
 &\le Q_n(R)D_*H^{D_*}\varepsilon
       +D_n(R)RH^{D_*+1}\varepsilon.
\end{align*}
The prefactor is bounded by a fixed power $H^a$, with a fixed
multiplicative constant. Choose $v>a+3/2$ before the final coupling
threshold. The right side is then at most $H^{-3/2}$ eventually.
This is the actual healthy local-coefficient error required by the
cumulant comparison. The support estimates of the marked propagation
put every variable used here at distance at most $B-2$ from the
chosen base; no estimate outside the healthy stencil has been used.

The atlas connected polynomial $c'_I(X)$ is the finite logarithm
of its jointly evaluated Gaussian moments. The exact common-coordinate
and calibration construction, followed by the actual covariance bounds,
compares those moments with the free moments with error $H^{-3/2}$.
Section~\ref{sec:allarity-cumulant-comparison} consequently gives
\[
 |c'_I(X)-c_I^{\rm free}|+H^{-3/2}\le H^{-31/32}
\]
eventually. Adding the healthy local error proves the prescribed
small remainder against the negative free connected coefficient.
The coordinate $X$ is common to all subsets, so this step does not
mix separately chosen background realizations.

At the upper comparison level the same construction starts with
singletons only. The transport has zero length there and is the
identity. Hence its moment polynomial is
$\prod_i(1+z_ip_i)$ in the squarefree marker algebra, and its
finite logarithm is $\sum_i z_ip_i$. Every connected coefficient
with at least two markers vanishes. The singleton is exactly the
calibrated local section; its support, marker and masked value agree
with the terminal birth term. The algebraic reindexing therefore
adds zero higher coefficients and leaves the pushed density unchanged.
Shrinking the marker radius preserves regular and forced bounds,
since every marker monomial acquires a smaller nonnegative weight.
This identifies the top initialization used in the all-arity marked
propagation with the actual calibrated insertion.

The remaining inputs in applying this chapter are the actual mixed
compiler's per-step row and covariance comparison, the terminal paid
insertion bound and the common purchases. The argument proves the
healthy-field and polynomial comparison from those precise receipts;
it does not produce them by changing the selected history.
